\def\paperRoman{V}
\documentclass[10pt,a4paper]{article}
\usepackage[margin=22mm,top=23mm,bottom=23mm]{geometry}
\usepackage{fontspec}
\setmainfont{Latin Modern Roman}
\setsansfont{DejaVu Sans}
\usepackage{babel}
\babelprovide[import,main]{spanish}
\usepackage{amsmath,amssymb}
\usepackage{unicode-math}
\setmathfont{Latin Modern Math}
\usepackage{graphicx,booktabs,longtable,array,enumitem,xcolor,fancyhdr,xurl,hyperref,tikz,microtype}
\usetikzlibrary{arrows.meta,positioning}
\definecolor{navy}{HTML}{17375C}
\definecolor{teal}{HTML}{0F8F86}
\definecolor{gold}{HTML}{B28B3C}
\hypersetup{colorlinks=true,linkcolor=navy,urlcolor=teal,citecolor=navy,pdfauthor={Juan Manuel Goenaga Castro}}
\setlength{\parindent}{0pt}
\setlength{\parskip}{5.2pt}
\linespread{1.07}
\setlist{nosep,leftmargin=1.5em,topsep=4pt}
\setlength{\headheight}{15pt}
\pagestyle{fancy}
\fancyhf{}
\fancyhead[L]{\sffamily\footnotesize GoToGym Quantum Research}
\fancyhead[R]{\sffamily\footnotesize Serie QAF audiovisual / \paperRoman}
\fancyfoot[L]{\sffamily\scriptsize Prepublicación metodológica / 22 de octubre de 2026}
\fancyfoot[R]{\sffamily\footnotesize\thepage}
\renewcommand{\headrulewidth}{0.4pt}
\newcommand{\papertitle}[3]{\begin{center}\sffamily\small\textcolor{teal}{MANUSCRITO #1 / ARQUITECTURA AUDIOVISUAL QAF}\par\vspace{9pt}{\huge\bfseries\textcolor{navy}{#2}}\par\vspace{7pt}{\large #3}\par\vspace{11pt}{\normalsize\bfseries Juan Manuel Goenaga Castro}\par\small GoToGym Quantum Research / Bogotá D.C., Colombia\par\small 22 de octubre de 2026 / Versión 1.0\end{center}\vspace{4pt}}
\newcommand{\status}[1]{\noindent\colorbox{navy!7}{\parbox{0.966\linewidth}{\small\textbf{Naturaleza del trabajo.} #1}}\par\vspace{6pt}}
\newcommand{\keywords}[1]{\textbf{Palabras clave:} #1.}
\newcommand{\R}{\mathbb{R}}
\newcommand{\C}{\mathbb{C}}
\newcommand{\tr}{\operatorname{Tr}}
\newcommand{\diag}{\operatorname{diag}}
\newcommand{\argmin}{\operatorname*{arg\,min}}
\newcommand{\argmax}{\operatorname*{arg\,max}}
\newcommand{\IP}{\mathcal{I}}
\newcommand{\sectionpage}[1]{\clearpage\section{#1}}
\newcommand{\compacttable}{\small\renewcommand{\arraystretch}{1.22}}
\emergencystretch=2em

\begin{document}
\papertitle{V}{Validación, competencias y ejecución del sistema}{Gobernanza multidisciplinaria y programa experimental de producción experta}
\status{Protocolo de desarrollo y validación. Reúne 56 perfiles profesionales y controles de competencia; no certifica que un asistente sustituya su experiencia práctica o sus responsabilidades.}
\begin{abstract}
Se propone una organización de investigación y producción que convierte los cuatro manuscritos anteriores en entregables verificables. Se identifican 56 perfiles que cubren ciencia, ingeniería, marca, dirección, producción, postproducción, sonido, accesibilidad y evaluación. Cada perfil se define por responsabilidad, artefacto y criterio de revisión. La competencia se verifica mediante casos, reproducción y evaluación independiente, en lugar de declararse por adopción de un rol. El protocolo compara el sistema propuesto con alternativas clásicas y producción profesional bajo recursos equivalentes; separa calidad percibida, exactitud de producto y coste operacional. Se presentan los controles algebraicos ejecutados, la evidencia heredada del prototipo y un diseño de estudio con asignación aleatoria y análisis jerárquico. Se especifican hitos, contratos, causas de rechazo y un piloto de running. La contribución central es una ruta para transformar propiedad intelectual y capacidades de software en películas auditables, con criterios explícitos para avanzar, revisar o descartar una técnica.
\end{abstract}
\keywords{validación audiovisual, competencias, profesiones, experimentos, reproducibilidad, QAF, gobernanza}

\section{Qué significa competencia experta en este sistema}
La aspiración de calidad mundial requiere decisiones comparables a las de profesionales y un proceso que permita detectar errores. No se obtiene por afirmar que un modelo ``es experto'' en una lista de profesiones. En esta arquitectura, un perfil es un conjunto de tareas, conocimientos, límites y entregables. Una persona puede cubrir varios perfiles; una función de software puede asistir a varios. La responsabilidad de validar un artefacto no desaparece al automatizar su producción.

Se propone una escala de competencia de cuatro niveles: ejecución guiada, ejecución autónoma en casos conocidos, resolución de variaciones y revisión competente de terceros. El nivel no se asigna por una autoevaluación del agente. Debe incluir casos desconocidos, fallos sembrados, reproducción y revisión independiente. Para dominios sensibles se conserva la intervención de un profesional pertinente.

La lista siguiente es un mapa amplio de \textbf{56 perfiles base para esta arquitectura}; no pretende enumerar todas las especialidades posibles del cine o la ciencia. Las revisiones de localización, idiomas y formato pueden requerir perfiles adicionales. La integración de saberes se materializa en checklists, modelos y pruebas, no en una supuesta adquisición instantánea de experiencia humana.

\clearpage
\section{Mapa profesional y artefactos exigidos}
{\compacttable
\begin{longtable}{p{.05\linewidth}p{.25\linewidth}p{.56\linewidth}}
\toprule N.º & Perfil & Responsabilidad y evidencia de competencia \\\midrule\endfirsthead
\toprule N.º & Perfil & Responsabilidad y evidencia de competencia \\\midrule\endhead
1 & Dirección científica & Pregunta, hipótesis, límites y revisión de transferencia entre dominios.\\
2 & Matemática aplicada & Modelos bien definidos, unidades, hipótesis y pruebas de consistencia.\\
3 & Redes tensoriales & Codificación, gauge, enlace, RDM y errores frente a referencia pequeña.\\
4 & Optimización & Formulación restringida, certificados y comparadores bajo el mismo presupuesto.\\
5 & Estadística & Diseño, potencia, intervalos y análisis por unidades independientes.\\
6 & Aprendizaje automático & Datos, particiones, ajuste, generalización y ausencia de fuga.\\
7 & Visión computacional & Geometría, movimiento, máscaras y límites de estimación visual.\\
8 & Ingeniería de conocimiento & Grafo de evidencia, procedencia y vínculo entre claim y plano.\\
9 & Arquitectura de sistemas & Contratos, fronteras de servicio, fallos y operación reproducible.\\
10 & Ingeniería de pipeline & Grafo de tareas, caché, revisiones y publicación de artefactos.\\
11 & Ingeniería GPU & Perfilado, memoria pico, paralelismo y coste real de ejecución.\\
12 & QA de software & Pruebas significativas, casos adversos y regresiones de entrega.\\
13 & Producción ejecutiva & Alcance, recursos, cronograma y aceptación de resultados.\\
14 & Dirección creativa & Concepto coherente entre propósito, audiencia y acabado.\\
15 & Dirección cinematográfica & Puesta en escena, interpretación, cámara y significado del montaje.\\
16 & Estrategia de marca & Identidad, tono, tokens y revisión de consistencia.\\
17 & Estrategia de marketing & Etapa comercial, promesa, CTA y métricas con denominador.\\
18 & Investigación de audiencia & Necesidades, comprensión y sesgos de selección del público.\\
19 & Guion & Conflicto, mecanismo, transformación y locución medible.\\
20 & Comunicación científica & Precisión, incertidumbre y traducción sin distorsión de capacidades.\\
21 & Storyboard & Planos ejecutables, continuidad y relación entre imagen y voz.\\
22 & Dirección de arte & Lenguaje visual, materiales y coherencia del mundo representado.\\
23 & Diseño de producción & Espacios, objetos y recursos realizables para cada escena.\\
24 & Dirección de fotografía & Exposición, óptica, luz, textura y referencias verificables.\\
25 & Operación de cámara & Encuadre, foco, movimiento y continuidad de toma.\\
26 & Iluminación y gaffer & Montaje de luz seguro, consistencia y separación del sujeto.\\
27 & Producción de campo & Localizaciones, logística y disponibilidad de recursos.\\
28 & Continuidad de rodaje & Vestuario, objetos, ejes, gestos y estados entre planos.\\
29 & Casting & Selección de personas/voz por criterios y contexto de marca.\\
30 & Interpretación & Acción humana creíble y continuidad de actuación.\\
31 & Vestuario & Identidad de prenda, ajuste, textura y continuidad.\\
32 & Maquillaje y caracterización & Apariencia consistente con luz, cámara y personaje.\\
33 & Modelado 3D & Geometría, escala, topología y activos editables.\\
34 & Texturizado & Resolución, detalle, origen y coherencia entre superficies.\\
35 & Look development y shading & Materiales, respuesta de luz y parámetros físicamente consistentes.\\
36 & Rigging & Controles, deformación estable y límites anatómicos del personaje.\\
37 & Animación & Peso, intención, ritmo y movimiento sin artefactos.\\
38 & Efectos y simulación & Fenómenos, unidades, estabilidad y reproducibilidad de simulación.\\
39 & Composición VFX & Integración, máscaras, alfa, luz, grano y profundidad.\\
40 & Motion graphics & Diagramas, tipografía y transiciones que expliquen el mecanismo.\\
41 & Colorismo & Transformaciones, balance, continuidad y salida revisada.\\
42 & Edición y montaje & Ritmo, continuidad, exactitud de duración y revisión narrativa.\\
43 & Lingüística y fonética regional & Pronunciación, variedad del español y evaluación contextual de acento.\\
44 & Dirección de voz & Interpretación, respiraciones, énfasis y consistencia de identidad.\\
45 & Locución & Dicción y actuación vocal; tomas y derechos de uso documentados.\\
46 & Diseño sonoro & Ambiente, efectos y motivos que apoyen la historia.\\
47 & Composición y supervisión musical & Identidad musical, edición y procedencia de pistas.\\
48 & Mezcla y mastering de audio & Balance, dinámica, sonoridad, picos y escucha final.\\
49 & Accesibilidad & Subtítulos, legibilidad, contraste y equivalencia de información.\\
50 & Especialista del dominio explicado & Revisión del contenido de entrenamiento, nutrición o salud según alcance.\\
51 & Gestión de derechos y licencias & Permisos de imagen, voz, música, modelos y distribución.\\
52 & Mastering y entrega de video & Máster, codec, cuadros, sincronía y reproducción de exportaciones.\\
53 & Seguridad y privacidad & Separación de cuentas, datos sintéticos y acceso a fuentes.\\
54 & MLOps y granja de render & Versiones, colas, observabilidad, reintentos e infraestructura.\\
55 & Economía de producción & Coste total, revisión, reutilización y análisis calidad--presupuesto.\\
56 & Investigación con participantes & Información al participante, protocolo, registro y custodia de evaluaciones.\\\bottomrule
\end{longtable}}

Un artefacto puede requerir firmas de varias competencias. El perfil de comunicación científica revisa el significado; el de fotografía revisa la imagen; el de dominio revisa el alcance especializado. Una nota estética no compensa una afirmación errónea. Tampoco una prueba matemática acredita una mezcla musical de calidad.

\section{Gobernanza de decisiones y contratos}
Cada módulo recibe una tarea acotada, evidencia y límites. Produce propuesta, incertidumbre, dependencias, coste y pruebas. El orquestador coordina estados y revisiones; no decide por autoridad retórica de un agente. Las decisiones críticas conservan una regla de abstención y un motivo verificable.

La compuerta de entrega es conjuntiva:
\begin{equation}
 G_{\mathrm{entrega}}=G_E\land G_D\land G_B\land G_A\land G_T\land G_R,
\end{equation}
para evidencia, derechos, marca, accesibilidad, controles técnicos y reproducción final. Los valores son estados revisados con procedencia. Un agente que escriba \texttt{true} sin realizar la comprobación no crea una certificación. Las revisiones independientes deben recibir el artefacto y la tarea, sin revelar la respuesta esperada cuando se pretende evaluar competencia.

Un error de alta severidad invalida el segmento dependiente y activa nueva versión. Se define una pérdida jerárquica: primero se minimizan violaciones críticas; después se optimizan comprensión, estética y recursos sobre el conjunto admisible. Esto evita que una suma ponderada permita ``pagar'' una falsedad con una mejora visual.

La protección de propiedad intelectual comienza con procedencia y acceso, no con una declaración automática de titularidad. Los materiales se clasifican como aportados, creados, licenciados o derivados, con permisos de distribución. Los documentos de esta serie citan a Antoni Bancells donde corresponde y conservan bibliografía de los métodos de terceros. No se realizó una búsqueda exhaustiva de patentes o libertad de operación.

\section{Resultados verificados y límites de esta serie}
Se ejecutó \texttt{validate\_models.py} con Python 3.12.14 y NumPy 2.3.5. El código y los resultados se incluyen en el paquete. Las comprobaciones son pequeñas y algebraicas:

\begin{table}[ht]\centering\compacttable
\begin{tabular}{p{.25\linewidth}p{.30\linewidth}p{.34\linewidth}}\toprule
Control & Resultado observado & Qué acredita \\\midrule
Cuatro estados de fase & Una diagonal; cuatro características complejas & Identificabilidad en un generador construido.\\
Aproximación de Bell & Fidelidad máxima de rango uno 0,5 & Necesidad de rango dos en ese ejemplo.\\
QUBO de seis variables & 64 estados, ocho factibles; mismo óptimo & Equivalencia de la compilación específica.\\
Cámara de mínimo jerk & Extremos correctos; integral 720,000072 & Concordancia numérica con 720 a tolerancia fijada.\\
Tiempo editorial & 1.800 cuadros a 30 fps & Duración de 60 s del esquema de cuadros.\\\bottomrule
\end{tabular}\caption{Controles nuevos ejecutados; no son un benchmark de espectadores o render.}
\end{table}

En el prototipo Studio v2 se habían ejecutado nueve pruebas del compositor y nueve escenarios en ocho motores de producto, además de un animatic gráfico de 60 s. Es evidencia antecedente del trabajo local; no una validación nueva de los modelos neuronales o del renderer físico aquí propuestos. La presente serie no ha producido una película final con rodaje, casting regional y render 3D integrado.

Las métricas de Q3--Q6 se citan como evidencia aportada por el programa QAF y con sus denominadores. No se reejecutaron sus bancos completos en esta tarea. El material de referencia identifica también regiones clásicas favorables, evidencia insuficiente y colapso computacional. El sistema conserva esas posibilidades como resultados válidos de la investigación.

\section{Diseño experimental de calidad y utilidad}
Se propone un estudio en dos etapas. El piloto usa 12 episodios y tres métodos: plantilla reproducible, planificación clásica profesional y planificación QAF. Se preparan 36 versiones verticales comparables. Un diseño equilibrado puede asignar a 60 personas seis episodios distintos por persona, dos por método, generando 360 evaluaciones y diez por versión. Nadie ve dos versiones del mismo episodio para reducir aprendizaje y comparación directa. Estos números son un diseño inicial, no una potencia estadística acreditada.

El desenlace primario es comprensión de mecanismo, con preguntas preespecificadas. Se registran confusión sobre capacidades, recuerdo, credibilidad, estética, autenticidad vocal y siguiente paso. El coste incluye producción, generación descartada y revisión. Se anotan fallos y no solo la mejor toma. El efecto del método puede estimarse con un modelo jerárquico:
\begin{equation}
 y_{epr}=\mu+\beta_{\mathrm{método}(r)}+u_e+v_p+\epsilon_{epr},
\end{equation}
donde $e$ identifica episodio y $p$ persona. Para respuestas binarias se usa enlace logístico; para escalas ordinales, un modelo ordinal. Los intervalos deben respetar agrupación por episodio y persona. Reutilizar muchas valoraciones de un único video no crea muchos videos independientes.

Para planear una comparación pareada confirmatoria, una aproximación de potencia es
\begin{equation}
 n\simeq\left[\frac{(z_{1-\alpha/2}+z_{1-\beta})\sigma_d}{\delta}\right]^2.
\end{equation}
Con $\alpha=0,05$, potencia 0,80, desviación de diferencias hipotética $\sigma_d=0,8$ y diferencia relevante $\delta=0,4$, resulta $n\approx32$ unidades pareadas independientes. Es un ejemplo de cálculo, no justifica tratar el piloto de 12 episodios como confirmatorio. El tamaño final se recalcula con varianzas y correlaciones del piloto.

Los desenlaces secundarios se preespecifican y se controla multiplicidad. Los pesos del score de calidad se fijan antes de abrir resultados. Se reporta el efecto y su incertidumbre, no únicamente un valor $p$. La mejora artística sin reducción de confusión puede ser una ganancia estética y un fracaso de comunicación.

\section{Baselines, ablaciones y regiones de no uso}
Los comparadores reciben las mismas fuentes, funciones disponibles, restricciones, datos de audiencia, presupuesto y número de intentos. Se iguala también el acceso a características relacionales complejas: compararlo solo con observaciones marginales no aísla la contribución de MPS o TDVP. Se utilizan modelos clásicos con grafo, secuencia y optimización cuando sean pertinentes.

Las ablaciones retiran, una por vez, relaciones, fase, dinámica, estado persistente y optimización; conservan lo demás. La comparación MPS contra estado denso en problemas pequeños mide fidelidad numérica, no superioridad sobre un algoritmo clásico. El contraste con la misma característica compleja y el mismo solver separa calidad de representación y efecto del optimizador.

Se prueban cambios de idioma, producto, audiencia, duración, presupuesto, escenas adversas y fuentes contradictorias. QAF se rechaza operacionalmente si la mejora no supera un umbral útil, desaparece fuera del generador o excede recursos. La categoría ``no identificable'' indica necesidad de mejores observaciones. No debe convertirse en una invitación a aumentar el tamaño del modelo sin datos adicionales.

\section{Programa de implementación y recursos}
\begin{table}[ht]\centering\compacttable
\begin{tabular}{p{.15\linewidth}p{.36\linewidth}p{.37\linewidth}}\toprule
Fase & Entregable & Condición para avanzar \\\midrule
1 & Fuentes, claims y contrato de episodio & Trazabilidad y alcance revisados.\\
2 & Guion, storyboard y audición & Comprensión preliminar y voz seleccionada.\\
3 & Pipeline clásico y máster & Render, color, mezcla y revisión reproducibles.\\
4 & Capa relacional QAF y ablaciones & Paridad de inputs, calidad y presupuesto.\\
5 & Piloto con audiencia y revisión & Efectos, fallos e incertidumbre publicados internamente.\\
6 & Serie y monitorización & Calidad sostenida y control de versiones.\\\bottomrule
\end{tabular}\caption{Hitos de trabajo; la duración se estima tras comprobar equipo, activos y acceso.}
\end{table}

El coste total se expresa como $C=\sum_r h_rv_r+C_{\mathrm{cómputo}}+C_{\mathrm{activos}}+C_{\mathrm{revisión}}$, con horas por rol $h_r$ y tarifa $v_r$ confirmada. No se asignan precios de mercado o una granja GPU ficticia. Un equipo pequeño puede cubrir varios perfiles; las revisiones críticas deben mantener independencia suficiente y competencia pertinente.

Los módulos iniciales son: gestor de fuentes, compilador de claims, timeline canónico, adaptadores de render, gestor de color, voz/alineación, mezcla, QA y registro experimental. La capa QAF se añade después de disponer de una referencia clásica medible. Cada módulo tiene interfaz, dataset de pruebas, versión y responsable. El primer entregable es una película revisable, no una plataforma enorme sin piloto terminado.

\section{Piloto de referencia y criterio de finalización}
El caso de running permite conectar situación humana, datos declarados y propuesta sin recurrir a capacidades futuras. Se requiere validar el recorrido real de la experiencia, seleccionar voz y producir tomas. El ejemplo sintético se rotula hasta disponer de captura autenticada y anonimizada. La narrativa no presenta la caminata como tratamiento ni como sesión ya realizada.

La película se considera lista al completar máster, versión de distribución, subtítulos, fuentes, proyecto editable e informe de revisión. Duración y codec son necesarios pero insuficientes. La revisión final debe reproducir imagen y sonido, comprobar pantalla móvil y confirmar que el público puede explicar la capacidad sin atribuciones falsas. La publicación no se deduce de la existencia de un archivo exportado.

\section{Conclusión}
La profesionalización del sistema exige integrar competencias con entregables, límites y pruebas. Esta serie aporta formalismo, dirección, arquitectura de render, identidad sonora y validación. El resultado actual es un programa científico y técnico ejecutable por etapas; la pretensión de ser la mejor tecnología del mundo permanece como una meta que necesita comparación independiente y evidencia sostenida.

\begin{thebibliography}{Q6}
\bibitem[Q1]{q1} Goenaga, J. M. (2025). \emph{Quantum Analysis Frameworks}. Aportado.
\bibitem[Q2]{q2} Bancells, A. (2026). \emph{QAF: dinámica operacional no conmutativa y compilación QUBO asíncrona}. Aportado.
\bibitem[Q3]{q3} Goenaga, J. M. (2026). \emph{Validación preliminar del QAF mediante benchmarking computacional reproducible}. Aportado.
\bibitem[Q4]{q4} Goenaga, J. M. (2026). \emph{Validación end-to-end del motor QAF: MPS + TDVP + 1-RDM/2-RDM}. Aportado.
\bibitem[Q5]{q5} Goenaga, J. M. (2026). \emph{Consolidated Experimental Evidence for QAF}. Aportado.
\bibitem[Q6]{q6} Goenaga, J. M. (2026). \emph{Value and Non-Value Regions of QAF}. Aportado.
\bibitem[R1]{rep} Goenaga, J. M. (2026). \emph{Controles algebraicos de QAF audiovisual}. Código, resultados, versiones y hashes adjuntos a esta serie.
\end{thebibliography}
\end{document}
